\documentclass[11pt,a4paper]{article}
\usepackage[margin=2.4cm]{geometry}
\usepackage{amsmath,amssymb}
\usepackage[colorlinks=true,linkcolor=blue,urlcolor=blue]{hyperref}
\usepackage{xcolor}
\usepackage{parskip}

\newenvironment{reviewercomment}%
{\par\medskip\noindent\begin{quotation}\itshape\color{black!70}}%
{\end{quotation}\medskip}
\newcommand{\response}{\par\noindent\textbf{Response:}\ }
\newcommand{\changes}{\par\noindent\textbf{Changes made:}\ }

\begin{document}

\begin{center}
{\Large\bfseries Response to Reviewer \#4}\\[0.4em]
{\large Ms.\ Ref.\ No.: MEAS-D-26-10707}\\[0.3em]
Gradient-Sensitive Functional Descriptors for Interferometric Surface Metrology:\\
Complementing Amplitude-Averaging ISO Parameters\\[0.3em]
\textit{Measurement} --- Revised version, September 2026
\end{center}

\bigskip

I thank Reviewer~\#4 for the demanding and detailed review. The four points call, in their fullest form, for a new experimental programme --- independent runs, thermal cycles, re-centring, a complete GUM Type-B budget, physical defects, additional materials and diameters, and an inter-laboratory comparison. This is a large scope, and not all of it can be executed within a single revision cycle by a single author. I have therefore followed the reviewer's own alternative instructions: I have done what the existing data genuinely support (a quantified type-B budget, a defect-to-background scaling analysis, and a standardised sampling/angular-unit treatment), and for what requires new experiments I have sharply narrowed the claims and defined a concrete validation programme rather than overstate the evidence. Each point is addressed below; a marked (redline) version of the manuscript accompanies this letter.

\subsection*{Point 4.1 --- Between-run reproducibility and a GUM-compliant uncertainty budget}
\begin{reviewercomment}
Expand the uncertainty analysis significantly throughout the manuscript. The current version relies solely on within run repeatability derived from five continuous sweeps [\ldots]. To address this major weakness you must provide a comprehensive between run reproducibility assessment. You are specifically required to include new experimental data involving completely independent measurement setups distinct thermal cycles and physical recentering of the measurement object. Furthermore the manuscript currently lacks a full GUM compliant uncertainty budget [\ldots]. These additions must include an assessment of laser wavelength stability refractive index variations Abbe error camera nonlinearity and form removal uncertainty.
\end{reviewercomment}
\response I have expanded the uncertainty analysis as far as the existing data permit, and I am explicit about what it cannot yet establish. (1)~A new subsection (Sec.~3.8, ``Indicative Type-B Uncertainty Budget'', with new Table~2) quantifies exactly the five contributions named: \emph{laser wavelength stability} ($\pm1$~nm module tolerance, rectangular: $0.11\%$ relative, $0.3$~nm on the FN~111 $RONt$; a pure scale factor that cannot create or mask localised features); \emph{refractive-index variation of air} ($<10^{-5}$ relative at the observed $\pm0.2^{\circ}$C stability: negligible); \emph{Abbe error / axis wobble} ($\pm12~\mu$rad, analysed through both the cosine-error and beam-walk mechanisms: sub-$0.01$~nm after LSCI first-harmonic removal); \emph{camera nonlinearity / EFM localisation} (bounded empirically by the $\pm0.05$-fringe per-frame uncertainty, reduced by band-limited averaging to $\approx1.1$~nm); and \emph{form-removal uncertainty} (lima\c{c}on second-order term $a^2/2R$, evaluated at $7$--$41$~nm for centring eccentricities of $15$--$35~\mu$m, carried as a bounded systematic because it is common to all sweeps, and consistent with the $34$~nm optical--tactile $RONt$ difference that the $E_n=0.61$ comparison shows to be within combined uncertainties). The combined random budget confirms the type-A term dominates and identifies the form-removal systematic and the untested between-run effects as the leading unquantified risks. (2)~Regarding new between-run experiments with independent setups, thermal cycles and physical re-centring: I agree these are required before the descriptors can be recommended for adoption, but they cannot be performed within this revision cycle, as the interferometer is no longer in the exact configuration used for the archives. Rather than present hastily acquired data, I state this limitation explicitly and in the reviewer's own terms. The paper's claims are scoped throughout as proof-of-concept and within-run only, and I believe they are commensurate with the evidence presented.
\changes New Sec.~3.8 and Table~2; Sec.~3.7 renamed ``Uncertainty Evaluation'' with a pointer to the type-B budget and the GUM reference (JCGM~100); Abstract, Limitations (first point) and Future Work item~(i) updated. I do not claim that five continuous sweeps provide full reproducibility; the text repeatedly states these are within-run lower bounds.

\subsection*{Point 4.2 --- Synthetic defects and translation to real-world conditions}
\begin{reviewercomment}
Improve and broaden the synthetic defect injection methodology to reflect real world conditions. [\ldots] If physical testing is absolutely impossible you need to substantially expand the discussion section. You must explain exactly how your synthetic sensitivity results can be reliably translated to practical industrial inspection scenarios such as the rigorous quality control of bearing surfaces complex optical components or precision mechanical seals.
\end{reviewercomment}
\response Calibrated physical-defect artefacts are not available to me within the revision period, so I have followed the reviewer's alternative instruction and substantially expanded the Discussion with a quantitative translation framework. A new paragraph in Sec.~5.3 derives a scaling rule: a localised defect adds a total-variation increment $\Delta TV$ that depends only on the defect geometry ($\Delta TV\approx 2A$ for a spike of amplitude $A$), so the relative full-profile BV-density response is $\Delta TV/TV_{\mathrm{background}}$. This reproduces the observed $+18\%$ on the FN~111 ($TV=1.137~\mu$m) and predicts only $+0.17\%$ for the identical spike on the FN~101 ($TV=118.4~\mu$m) --- a direct quantification of the reviewer's observation that the FN~101 defect-to-background ratio is roughly 500 times less favourable. Two practical consequences are drawn: full-profile BV density is best suited to smooth finished surfaces (bearing raceways, optical components, precision seals) whose residual amplitudes are comparable to the FN~111 case, with the synthetic archetypes mapped to realistic feature classes (spike~$\leftrightarrow$~pit/inclusion edge, scratch~$\leftrightarrow$~tooling or handling mark, sinusoid~$\leftrightarrow$~chatter); and for rougher backgrounds the sliding-window formulation restores local contrast, which is exactly why the Flick flat is detected on the FN~101 despite the unfavourable full-profile ratio. The text states plainly that validation against physical defects of calibrated dimensions remains necessary before these scaling arguments can be considered established.
\changes New paragraph in Sec.~5.3; Limitations third point rewritten to reference the scaling analysis; Future Work item~(iii) upgraded to calibrated physical defects and a blind-detection protocol.

\subsection*{Point 4.3 --- Sampling-interval dependence and angular-unit conventions}
\begin{reviewercomment}
Address the complex sampling interval dependence of the newly proposed functional descriptors. [\ldots] they require a rigorous and mathematically sound conversion to angular units such as nanometers per degree for any meaningful cross instrument comparison. [\ldots] You must provide a clear standardized procedure [\ldots]. The development and proposal of agreed angular unit conventions is absolutely critical at this stage.
\end{reviewercomment}
\response I agree this is the most technically important point, and I have turned it into an explicit element of the methodology (new Sec.~3.6, ``Angular-Unit Convention for Sampling-Dependent Descriptors''). The convention has three steps: (1)~\emph{band-limit first} --- apply a declared phase-correct Gaussian profile filter (ISO~16610-21) at a stated UPR cutoff, after which the profile derivative is bounded and the discrete sums converge to the continuous integrals as $\Delta\alpha\to0$, so the descriptor characterises the declared bandwidth rather than the sampling grid; (2)~\emph{sample adequately} --- at least ten samples per cutoff undulation (the present data carry $\approx$290 and $\approx$1200); (3)~\emph{report in angular and arc-length units together with the filter cutoff}. Crucially, I now make the dimensional distinction the reviewer's comment points to: amplitude descriptors ($R_a$, $R_q$, $L^1$, $L^2$, $RONt$) have the dimension of length (nm) and are sampling-independent, whereas the gradient descriptors are height-per-length and must be reported as nm\,deg$^{-1}$ and, for diameter-independent cross-instrument comparison, as nm\,mm$^{-1}$ of arc length ($\mathrm{BV}_s=TV/\pi D$). Converted values are given for both artefacts (FN~111: BV $3.2$~nm\,deg$^{-1}=12.1$~nm\,mm$^{-1}$; FN~101: BV $329$~nm\,deg$^{-1}=1.26~\mu$m\,mm$^{-1}$, at the 15-UPR cutoff), and the text notes that the nearly identical circumferences ($\approx94$~mm) make the large arc-length difference a genuine surface effect rather than a unit artefact. Because each conversion is a fixed multiplicative factor, relative uncertainties and relative defect responses are identical in every unit system. The parallel with declaring the filter for $RONt$ is drawn explicitly as the route to reference-software adoption; cross-instrument empirical verification is listed as Future Work.
\changes New Sec.~3.6 (with the arc-length nm\,mm$^{-1}$ convention and the amplitude-versus-gradient dimensional distinction); Sec.~3.5 closing paragraph updated to point to it; \emph{Practical considerations} (Sec.~5.3), Limitations (fourth point), Metrological Implications and Future Work updated to reference the convention.

\subsection*{Point 4.4 --- Validation scope and generalizability}
\begin{reviewercomment}
Broaden the validation scope and definitively demonstrate the generalizability of your findings. The current experimental validation is far too narrow [\ldots] strictly limited to merely two specific artifacts from a single commercial manufacturer JENOPTIK. [\ldots] You need to rigorously validate your proposed methodology on a significantly wider range of artifact types [\ldots]. Finally conducting detailed inter laboratory comparisons using a common artifact set would provide the vital long term reproducibility data [\ldots].
\end{reviewercomment}
\response I agree that generalisability is not yet demonstrated, and the revision now says so without qualification. New certified artefacts of different diameters and materials (optical glass, tungsten carbide) and an inter-laboratory campaign are beyond what a single author can add in this revision cycle; they require procurement and partner laboratories. The revision therefore (a)~scopes every conclusion explicitly as a proof-of-concept on two artefacts; (b)~notes the breadth the existing data do provide --- the two artefacts differ substantially in material (Al$_2$O$_3$ ceramic versus steel) and by two orders of magnitude in residual amplitude, though they share diameter and manufacturer; and (c)~converts the reviewer's requirements into a concrete validation programme in Future Work: artefacts of different diameters, materials and finishes; calibrated physical defects and blind detection; an inter-laboratory comparison with a common artefact set and standardised pipeline as the source of long-term reproducibility data for potential standardisation; and cross-instrument verification of the angular-unit convention. I believe the revised claims are now strictly commensurate with the validation scope.
\changes Limitations second point rewritten; Future Work items~(ii)--(v) expanded; proof-of-concept scoping verified throughout (Abstract, Sec.~5.3, Sec.~5.5).

\bigskip
I hope the reviewer finds that the revision responds substantively to every point, within the constraints of what a single author can accomplish by the resubmission deadline, and that the claims are now properly bounded by the evidence.

\bigskip
Yours sincerely,\\[0.6em]
Dawid Kucharski\\
Institute of Mechanical Technology, Poznan University of Technology\\
\href{mailto:dawid.kucharski@put.poznan.pl}{dawid.kucharski@put.poznan.pl}

\end{document}
